\chapter{Limping}

\section*{Definition}
Limping (antalgic gait) is an abnormal gait pattern that minimizes weight-bearing on a painful lower extremity, characterized by a shortened stance phase on the affected side and a compensatory prolonged swing phase or trunk lean.

\section*{What Happens in Your Body}
Nociceptive signals from the hip, knee, ankle, or foot activate protective muscle splinting and reduce loading forces via reciprocal inhibition of ipsilateral extensors. Gait cycle analysis shows decreased single-limb stance time on the painful side, often with hip drop (Trendelenburg sign) or vaulting (toe-walking) to clear the painful limb during swing.

\section*{Causes}
\begin{itemize}
  \item \textbf{Traumatic:} Stress fracture, acute fracture, sprain (lateral ankle most common), ligamentous tear (ACL, MCL), meniscal tear, contusion, foot or heel pain (plantar fasciitis, Achilles tendinopathy)
  \item \textbf{Arthritic:} Osteoarthritis (hip or knee), juvenile idiopathic arthritis, rheumatoid arthritis, gout, pseudogout, septic arthritis (emergent)
  \item \textbf{Pediatric:} Transient synovitis of the hip (most common cause age 3--10), Legg--Calv\'{e}--Perthes disease (age 4--8), slipped capital femoral epiphysis (SCFE, age 10--16), developmental dysplasia of the hip
  \item \textbf{Infectious:} Septic arthritis, osteomyelitis, Lyme disease, soft tissue abscess
  \item \textbf{Vascular:} Avascular necrosis (femoral head), chronic venous insufficiency with ulceration, peripheral artery disease (claudication)
  \item \textbf{Neuromuscular:} Cerebral palsy (spastic hemiplegia), stroke with hemiparesis, hereditary spastic paraplegia, Charcot--Marie--Tooth disease (foot drop), myopathy
  \item \textbf{Miscellaneous:} Ilio-tibial band syndrome, patellofemoral pain syndrome, leg-length discrepancy, inguinal hernia, testicular torsion (referred pain)
\end{itemize}

\section*{When to See a Doctor}
See your doctor if:
\begin{itemize}
  \item Limp is acute and you cannot bear any weight (non-weight-bearing emergency)
  \item Fever, erythema, warmth, or swelling of a joint (consider septic arthritis until proven otherwise)
  \item Limp persists for \textgreater{}2 weeks without clear traumatic cause
  \item Child refuses to walk or crawl (requires urgent evaluation)
  \item Night pain, systemic symptoms (unexplained fever, night sweats, weight loss), or history of malignancy
  \item Neurologic signs: foot drop, sensory loss, spasticity, or asymmetric reflexes
\end{itemize}

\section*{Related Symptoms}
\begin{itemize}
  \item Joint pain (hip, knee, ankle), swelling or effusion, stiffness (morning or after rest), reduced range of motion, muscle atrophy, Trendelenburg gait, short-leg gait, foot drop (steppage gait), hip or groin pain (referred from hip joint)
\end{itemize}
\textbf{Important:} In children, the age of presentation narrows the differential considerably: transient synovitis is common at age 3--10, SCFE at age 10--16, and Perthes at age 4--8; never assume a limp is benign without evaluating for these conditions.

\section*{Self-Care / Home Management}
\begin{itemize}
  \item Rest the affected limb, avoiding weight-bearing activities that reproduce pain
  \item Apply ice to the painful area for 15--20 minutes every 2--3 hours if acute injury
  \item Use crutches or a cane (held in the contralateral hand) to unload the painful joint
  \item Take over-the-counter analgesics (acetaminophen or NSAIDs like ibuprofen) as tolerated if no contraindications
  \item Avoid prolonged sitting or standing; gentle stretching of the unaffected side
\end{itemize}

\section*{How Your Doctor Will Evaluate You}
\begin{itemize}
  \item \textbf{History:} Onset (acute vs. chronic), trauma, fever, night pain, morning stiffness, referred pain patterns, age, activity level, recent infections
  \item \textbf{Physical exam:} Observation of gait (barefoot); inspection for swelling, erythema, leg-length discrepancy; palpation for point tenderness; range of motion of hip, knee, ankle; FABER and FADIR tests (hip); log roll test; neurological exam (strength, reflexes, sensation, tone)
  \item \textbf{Imaging:} Plain radiographs (AP pelvis and bilateral frog-leg lateral views for hip pathology; weight-bearing views of knee/ankle if indicated); ultrasound for effusion or joint fluid; MRI or CT if radiographs equivocal and concern for occult fracture, AVN, or osteomyelitis
  \item \textbf{Laboratory:} CBC, ESR, CRP; blood cultures and joint aspiration with Gram stain, culture, cell count, and crystal analysis if septic arthritis suspected; Lyme serology in endemic areas
\end{itemize}

\section*{Treatment Options}
\begin{itemize}
  \item \textbf{Traumatic:} RICE (rest, ice, compression, elevation) for acute injuries; NSAIDs; physical therapy; bracing or casting as indicated; orthopedic referral for fractures, complete ligament tears, or meniscal injuries
  \item \textbf{Arthritic:} NSAIDs, intra-articular corticosteroid injections, weight loss, physical therapy, joint replacement for end-stage osteoarthritis
  \item \textbf{Pediatric (hip):} Transient synovitis managed with rest and NSAIDs; SCFE requires urgent surgical pinning; Perthes treated with containment (bracing or surgery) if diagnosed early
  \item \textbf{Infectious:} Septic arthritis requires emergent irrigation and debridement plus IV antibiotics; osteomyelitis treated with tailored IV antibiotics and possible surgical drainage
  \item \textbf{Neuromuscular:} Ankle--foot orthosis (AFO) for foot drop; nerve conduction studies if progressive; botulinum toxin for spasticity; referral to neurology or physiatry
  \item \textbf{Leg-length discrepancy:} Shoe lift (1/3 to 2/3 of discrepancy); consider epiphysiodesis in skeletally immature people if discrepancy is projected to exceed 2 cm
\end{itemize}

\section*{References}
\begin{thebibliography}{99}
\bibitem{limping:ref1} Leet AI, Skaggs DL. "Evaluation of the acutely limping child." \textit{American Family Physician}, 2000;61(7):2105--2114.
\bibitem{limping:ref2} Sawyer JR, Kapoor M. "The limping child: a systematic approach to diagnosis." \textit{American Family Physician}, 2009;79(3):215--224.
\bibitem{limping:ref3} Perry DC, Bruce C. "Evaluating the child who presents with an acute limp." \textit{British Medical Journal}, 2010;341:c4250.
\bibitem{limping:ref4} Kocher MS, Zurakowski D, Kasser JR. "Differentiating between septic arthritis and transient synovitis of the hip in children." \textit{Journal of Bone and Joint Surgery}, 1999;81(12):1662--1670.
\bibitem{limping:ref5} Novais EN. "Limping child: a challenge in primary care." \textit{Journal of Pediatric Orthopaedics}, 2020;40(Suppl 1):S13--S17.
\bibitem{limping:ref6} Fischer SU, Beattie TF. "The limping child: epidemiology and assessment." \textit{Pediatrics}, 1999;103(6):e79.

\end{thebibliography}
